\begin{theorem}[Conditional fixed-tape multiplication bound]
Assume the pinned upstream multiplication theorem and retained analytic,
ordered-affine, compact-control, Gaussian-map, and fixed-tape interfaces
specified in the opening section. The partial-swap interchange interface
and the complex layer proved here may be substituted into those
interfaces. They give
\[
 T(n)=O\!\left(n(\log n)^{1-\kappa}\right),\qquad
 \boxed{\kappa=\frac{942293}{500000000000}
                 =1.884586\times10^{-6}}.
\]
\end{theorem}
\begin{proof}
Take the following exact parameters:
\[
\begin{aligned}
 a_b&=\frac{188459}{50000000000},&a_c&=\frac{417}{100000000},\\
 \tau&=1-a_b,&\sigma&=1-a_c,\\
 \epsilon&=\frac1{2+a_b}=\frac{50000000000}{100000188459},&c&=1,\\
 \beta&=\frac1{1000},&\delta&=\frac1{10^{16}},\\
 \lambda&=\tau+\frac1{10^{20}},&\lambda'&=\tau+\frac2{10^{20}},\\
 \zeta&=\frac1{10000},&C_1&=\frac{11999}{10000}.
\end{aligned}
\]
The strict bit and complex moments were certified above. Since
$\sigma<\tau$, the internal layer exponent is $\chi=\tau$.
The complex leaf saving is
\[
 (1-\beta)a_c=\frac{416583}{10^{11}}=4.16583\times10^{-6}>a_b.
\]
Thus both the internal and leaf exponents are strictly below
$\lambda'$, and the displayed gaps between $\tau,\lambda,\lambda'$
absorb fixed logarithmic factors. The precision guard has
$C_1=6/5-\beta/5+\zeta$ and $1-\epsilon C_1>0$.
Compact control allows $c=1$; its reservation exponent is zero and
$1-\epsilon(1+c)=1-2\epsilon>0$ gives $K=o(\ell)$ with $K=d$.

For the retained Gaussian parameters
$\alpha=\lfloor\sqrt{b/(8d)}\rfloor$ and
$\gamma_G=2d\alpha^2\le b/4$, the same positive $1-2\epsilon$
ensures the eventual conditions $b\ge64d^2$, $\alpha^2\theta_i\ge1$,
and prime-interval growth. Use the corrected chirp input bound
$F=2^{\lceil1.14\alpha^2\rceil}$; the following section proves that
$P=34p$ still suffices. The retained target errors and final rounding
therefore apply.

The seven assembly savings are the exact rational numbers
\[
\begin{aligned}
 g_1&=1-2\epsilon,&g_2&=\epsilon a_b,&
 g_3&=\epsilon(a_b-2/10^{20}),\\
 g_4&=(1-\epsilon)a_b,&g_5&=1-\delta-2\epsilon,&
 g_6&=1-\delta-\epsilon,&g_7&=\epsilon.
\end{aligned}
\]
Writing $b_0=a_b/(2+a_b)$ gives $g_1=g_2=b_0$,
$g_3=b_0-2\epsilon/10^{20}$, $g_4>b_0$,
and $g_5=b_0-\delta$. Since $\delta>2\epsilon/10^{20}$,
$g_5$ is the least of the seven. Exact rational subtraction gives
\[
 \min_i g_i-\kappa
 =\frac{a_b}{2+a_b}-\frac1{10^{16}}
       -\frac{942293}{500000000000}
 >\frac{44822}{10^{17}}>0.
\]
All retained parameter conditions are strict with this choice; in
particular the finite arithmetic certificate evaluates the complete
29-condition interface and these seven margins with rational arithmetic.
The final positive gap absorbs remaining logarithmic factors, proving
the claimed bound in the retained model.
\end{proof}

The count and moment data are finite arithmetic certificates for the
specified producer and common-basis construction. The theorem's use of
the pinned analytic interfaces and its fixed rational basis preparation
is separate from checking those integers and rational inequalities.
